\chapter{The Only Fast Output}
% full version: chapters/13_muscles.tex

Everything you do quickly in the world---a word, a glance, a step---is a contraction of muscles. The machine also has a second, slow output, a chemical one; Chapter Sixteen is about it. Here we deal with muscles.

The last link in the chain is formed by the \nt{ACET} neurons. Each muscle fiber listens to exactly one such neuron, and one neuron controls a group of fibers, from hundreds to a couple of thousand. Interestingly, the junction between nerve and muscle is the most reliable connection in the body. Inside the brain, connections lose more than half of their signals, but here each click releases several times more substance than needed. At the output an error costs a movement, and reliability is bought with a margin.

\section*{A Drumroll Becomes a Push}

One click, one short twitch of the muscle. If the clicks come rarely, the muscle jerks in separate twitches. If they come often, the twitches merge into a steady force, the way a drumroll merges into a rumble. The muscle turns a stream of spikes into force, like the simplest digital-to-analog converter.

\pencilfig{113}{%
% twitch = alpha function; the sum of twitches for spikes 1 s and 0.16 s apart (arbitrary units of time)
\begin{scope}[x=0.95cm, y=0.36cm]
\foreach \dx/\t in {0/rare,4.6/frequent}{
  \begin{scope}[xshift=\dx cm]
  \draw[pen] (0,0) -- (4,0); \node[font=\small\itshape, text=pcGray, anchor=south] at (2,5.4) {\t};
  \end{scope}}
\draw[penlite=pcRed] plot[smooth] coordinates {(0.00,0.00) (0.05,0.00) (0.10,0.00) (0.15,0.00) (0.20,0.00) (0.25,0.45) (0.30,0.73) (0.35,0.90) (0.40,0.98) (0.45,1.00) (0.50,0.98) (0.55,0.94) (0.60,0.88) (0.65,0.81) (0.70,0.74) (0.75,0.66) (0.80,0.59) (0.85,0.52) (0.90,0.46) (0.95,0.41) (1.00,0.35) (1.05,0.31) (1.10,0.27) (1.15,0.23) (1.20,0.20) (1.25,0.62) (1.30,0.88) (1.35,1.02) (1.40,1.08) (1.45,1.09) (1.50,1.06) (1.55,1.00) (1.60,0.93) (1.65,0.86) (1.70,0.78) (1.75,0.70) (1.80,0.62) (1.85,0.55) (1.90,0.48) (1.95,0.42) (2.00,0.37) (2.05,0.32) (2.10,0.28) (2.15,0.24) (2.20,0.21) (2.25,0.62) (2.30,0.88) (2.35,1.03) (2.40,1.09) (2.45,1.09) (2.50,1.06) (2.55,1.01) (2.60,0.94) (2.65,0.86) (2.70,0.78) (2.75,0.70) (2.80,0.62) (2.85,0.55) (2.90,0.48) (2.95,0.42) (3.00,0.37) (3.05,0.32) (3.10,0.28) (3.15,0.24) (3.20,0.21) (3.25,0.62) (3.30,0.88) (3.35,1.03) (3.40,1.09) (3.45,1.09) (3.50,1.06) (3.55,1.01) (3.60,0.94) (3.65,0.86) (3.70,0.78) (3.75,0.70) (3.80,0.62) (3.85,0.55) (3.90,0.48) (3.95,0.42) (4.00,0.37)};
\foreach \s in {0.20,1.20,2.20,3.20}{\draw[pcBlue, line width=0.8pt] (\s,-0.35) -- (\s,-1.2);}
\begin{scope}[xshift=4.6cm]
\draw[penlite=pcRed] plot[smooth] coordinates {(0.00,0.00) (0.05,0.00) (0.10,0.00) (0.15,0.00) (0.20,0.00) (0.25,0.45) (0.30,0.73) (0.35,0.90) (0.40,1.35) (0.45,1.68) (0.50,1.85) (0.55,2.19) (0.60,2.51) (0.65,2.64) (0.70,2.84) (0.75,3.13) (0.80,3.22) (0.85,3.27) (0.90,3.55) (0.95,3.62) (1.00,3.54) (1.05,3.82) (1.10,3.88) (1.15,3.80) (1.20,3.98) (1.25,4.06) (1.30,3.97) (1.35,4.07) (1.40,4.16) (1.45,4.09) (1.50,4.10) (1.55,4.23) (1.60,4.17) (1.65,4.10) (1.70,4.26) (1.75,4.23) (1.80,4.07) (1.85,4.27) (1.90,4.27) (1.95,4.12) (2.00,4.26) (2.05,4.29) (2.10,4.17) (2.15,4.24) (2.20,4.31) (2.25,4.21) (2.30,4.21) (2.35,4.31) (2.40,4.24) (2.45,4.16) (2.50,4.31) (2.55,4.27) (2.60,4.10) (2.65,4.30) (2.70,4.29) (2.75,4.15) (2.80,4.28) (2.85,4.31) (2.90,4.18) (2.95,4.25) (3.00,4.32) (3.05,4.22) (3.10,4.21) (3.15,4.32) (3.20,4.25) (3.25,4.17) (3.30,4.31) (3.35,4.27) (3.40,4.11) (3.45,3.86) (3.50,3.56) (3.55,3.25) (3.60,2.93) (3.65,2.63) (3.70,2.33) (3.75,2.06) (3.80,1.81) (3.85,1.58) (3.90,1.38) (3.95,1.19) (4.00,1.03)};
\foreach \s in {0.20,0.36,0.52,0.68,0.84,1.00,1.16,1.32,1.48,1.64,1.80,1.96,2.12,2.28,2.44,2.60,2.76,2.92,3.08,3.24}{\draw[pcBlue, line width=0.8pt] (\s,-0.35) -- (\s,-1.2);}
\end{scope}
\node[lbl, text=pcRed, anchor=west] at (1.2,1.9) {twitches};
\node[lbl, text=pcRed, anchor=south] at (6.6,4.55) {steady force};
\node[lbl, text=pcBlue, anchor=north] at (4.3,-1.35) {the neuron's clicks};
\end{scope}
}{Rare clicks give separate twitches of the muscle; frequent ones merge into a steady and much larger force.}

\section*{Floating-Point Force}

The brain has two force knobs: how often the neurons click and how many groups of fibers are switched on. The groups switch on in order, from the weakest to the strongest, and nobody computes this order: small neurons are simply easier to excite. As a result, each new group adds roughly the same \emph{fraction} to the force already there. When you pick up an egg, the force is metered out in tiny portions; when you lift a suitcase, in large ones. An engineer will recognize a floating-point number here: the exponent is set by the number of groups switched on, and the exact value by the rate of clicks. The flip side: the stronger the movement, the less precise it is. According to one influential hypothesis, this is why our movements are so smooth: a smooth command gives the least scatter.

\section*{Pull, Don't Push}

A muscle can only pull. So each joint is served by a pair of muscles pulling in opposite directions: again two wires instead of a sign. The difference of their forces turns the joint, and their sum makes it stiffer. When you carry a full cup, you tense both muscles at once: the rotation is the same, but the arm takes a bump more firmly.

\pencilfig{13}{\pencilart{13}}{A muscle can only pull, so there are two: the difference of the forces turns the joint, the sum makes it stiffer.}

\section*{Feedback with a Delay}

Muscles have stretch sensors, and the shortest loop closes right in the spinal cord: stretch a muscle, and it contracts by itself, while the opposing muscle is switched off for that time. It is switched off by a braking neuron of the \nt{GLYC} type: so this type has found its job.

But every loop has a delay: tens of milliseconds in the spinal cord, hundreds through the eye. And delayed feedback makes a system swing, like a driver who steers while seeing the road with a delay. The longer the loop, the more slowly it must correct. The limit beyond which the spinal loop starts to swing is about eight oscillations per second. This is close to the frequency of the fine tremor that every healthy person has, and the stretch loop is one of its likely sources.

So how do we move fast and precisely? According to a widespread hypothesis, we predict: we keep an internal model of the body and compute the result of a command without waiting for the sensors.

\section*{Rhythm Without a Metronome}

Walking and breathing are rhythmic, but they need no metronome. Two groups of neurons that inhibit each other and gradually tire take turns by themselves: the winner tires, the rested rival takes over, and so on again and again. A constant command from above, ``walk,'' turns on the spot into a ``left, right'' rhythm.

\fullversion{13}
